\documentclass[border=10pt]{standalone}
\usepackage{tikz,ctex}
\usetikzlibrary{positioning}

\begin{document}

\begin{tikzpicture}[
    % 全局设置
    node distance = 1.6cm and 1.4cm,
    every node/.style = {
        draw,
        rounded corners,
        thick,
        align = center,
        font = \sffamily\small,
        minimum width = 2.0cm,
        minimum height = 0.9cm
    },
    % 样式定义
    centerStyle/.style = {
        fill = blue!45,
        text = white,
        font = \sffamily\bfseries\large,
        line width = 2pt,
        minimum size = 2.8cm,
        circle
    },
    branchTitle/.style = {
        font = \sffamily\bfseries,
        line width = 1.5pt
    },
    subNode/.style = {
        fill = white,
        font = \sffamily\small,
        minimum width = 2.2cm,
        minimum height = 0.8cm
    },
    leafNode/.style = {
        fill = white,
        font = \sffamily\footnotesize,
        minimum width = 1.9cm,
        minimum height = 0.7cm
    },
    arrowStyle/.style = {
        ->,
        thick,
        >=stealth,
        shorten >=2pt
    }
]

    % ================= 1. 中心节点 =================
    \node[centerStyle] (center) {9.1.归纳偏置与不变性 \\ Inductive Bias \& \\ Invariance};

    % ================= 2. 上方分支：逆问题与归纳偏置 (红) =================
    \node[branchTitle, fill=red!15, above=of center] (inverse) {逆问题与归纳偏置 \\ Inverse Problems \& Inductive Bias};

    \node[subNode, above=of inverse, xshift=-6cm] (illposed) {不适定逆问题 \\ Ill-posed Inverse};
    \node[subNode, above=of inverse, xshift=-2cm] (bias) {归纳偏置 \\ Inductive Bias};
    \node[subNode, above=of inverse, xshift=2cm] (prior) {先验知识 \\ Prior Knowledge};
    \node[subNode, above=of inverse, xshift=6cm] (transfer) {迁移与多任务 \\ Transfer \& Multi-task};

    % ================= 3. 右侧分支：对称性与不变性 (蓝) =================
    \node[branchTitle, fill=orange!15, right=of center] (symmetry) {对称性与不变性 \\ Symmetry \& Invariance};

    \node[subNode, right=of symmetry, yshift=4cm] (trans) {平移不变性 \\ Translation Invariance};
    \node[subNode, right=of symmetry, yshift=2cm] (scale) {尺度不变性 \\ Scale Invariance};
    \node[subNode, right=of symmetry, yshift=0cm] (group) {群论与四公理 \\ Group Theory};
    \node[subNode, right=of symmetry, yshift=-2cm] (geometric) {几何深度学习 \\ Geometric DL};
    \node[subNode, right=of symmetry, yshift=-4cm] (learnhard) {直接学习的困难 \\ Hard to Learn Directly};

    % ================= 4. 下方分支：等变性 (绿) =================
    \node[branchTitle, fill=green!15, below=of center] (equivar) {等变性 \\ Equivariance};

    \node[subNode, below=of equivar, xshift=-6cm] (def) {等变定义 \\ $\mathcal{S}(\mathcal{T}(I))=\mathcal{T}(\mathcal{S}(I))$};
    \node[subNode, below=of equivar, xshift=-2cm] (seg) {图像分割例子 \\ Segmentation};
    \node[subNode, below=of equivar, xshift=2cm] (general) {广义等变 \\ $\tilde{\mathcal{T}}$ 不同变换};
    \node[subNode, below=of equivar, xshift=6cm] (special) {不变性是特例 \\ Invariance as Special Case};

    % ================= 5. 左侧分支：没有免费午餐与实现途径 (紫) =================
    \node[branchTitle, fill=purple!15, left=of center] (nfl) {没有免费午餐与实现途径 \\ No Free Lunch \& Approaches};

    \node[subNode, left=of nfl, yshift=5cm] (nfltheorem) {没有免费午餐定理 \\ No Free Lunch};
    \node[subNode, left=of nfl, yshift=3cm] (modelbased) {基于模型的机器学习 \\ Model-based ML};
    \node[subNode, left=of nfl, yshift=1cm] (pre) {预处理 \\ Pre-processing};
    \node[subNode, left=of nfl, yshift=-1cm] (reg) {正则化误差 \\ Regularized Error};
    \node[subNode, left=of nfl, yshift=-3cm] (aug) {数据增强 \\ Data Augmentation};
    \node[subNode, left=of nfl, yshift=-5cm] (arch) {网络结构 \\ Architecture};

    % ================= 连线逻辑 =================
    % 中心 -> 四大分支标题
    \draw[arrowStyle, red] (center) -- (inverse);
    \draw[arrowStyle, blue] (center) -- (symmetry);
    \draw[arrowStyle, green!60!black] (center) -- (equivar);
    \draw[arrowStyle, purple] (center) -- (nfl);

    % 分支标题 -> 子节点 (上方)
    \draw[arrowStyle, red!60] (inverse) -- (illposed.south);
    \draw[arrowStyle, red!60] (inverse) -- (bias.south);
    \draw[arrowStyle, red!60] (inverse) -- (prior.south);
    \draw[arrowStyle, red!60] (inverse) -- (transfer.south);

    % 分支标题 -> 子节点 (右侧)
    \draw[arrowStyle, blue!60] (symmetry) -- (trans.west);
    \draw[arrowStyle, blue!60] (symmetry) -- (scale.west);
    \draw[arrowStyle, blue!60] (symmetry) -- (group.west);
    \draw[arrowStyle, blue!60] (symmetry) -- (geometric.west);
    \draw[arrowStyle, blue!60] (symmetry) -- (learnhard.west);

    % 分支标题 -> 子节点 (下方)
    \draw[arrowStyle, green!60!black] (equivar) -- (def.north);
    \draw[arrowStyle, green!60!black] (equivar) -- (seg.north);
    \draw[arrowStyle, green!60!black] (equivar) -- (general.north);
    \draw[arrowStyle, green!60!black] (equivar) -- (special.north);

    % 分支标题 -> 子节点 (左侧)
    \draw[arrowStyle, purple!60] (nfl) -- (nfltheorem.east);
    \draw[arrowStyle, purple!60] (nfl) -- (modelbased.east);
    \draw[arrowStyle, purple!60] (nfl) -- (pre.east);
    \draw[arrowStyle, purple!60] (nfl) -- (reg.east);
    \draw[arrowStyle, purple!60] (nfl) -- (aug.east);
    \draw[arrowStyle, purple!60] (nfl) -- (arch.east);

\end{tikzpicture}

\end{document}